\begin{abstract}
Low-precision rollouts can make reinforcement learning (RL) of large language models cheaper, but they sample from a quantized copy of the policy being trained. QuRL \citep{li2026qurl} argues that re-quantizing the policy at every step masks the small RL weight updates. It proposes Update-Aware Quantization (UAQ): a one-time rescaling of the query, key and value weights by $1/s$, absorbed into the preceding normalization gain. According to QuRL, this reduces quantization error $s$-fold while amplifying updates by the same factor.

We show that the error reduction does not exist and that the amplification is a change of optimizer settings.
\begin{itemize}[leftmargin=*,itemsep=1pt,topsep=2pt]
\item \textbf{No error reduction.} The rescaling leaves the quantized forward pass unchanged for every quantizer whose scales are computed from absolute maxima, including QuRL's INT8 scheme (per-channel weights, per-token activations). This is exact whenever $s$ is a power of two.
\item \textbf{An optimizer reparameterization.} Under AdamW, training in UAQ coordinates equals training the original model with the learning rate multiplied by $s$ on the query/key/value weights and divided by $s$ on the normalization gains, with Adam's $\varepsilon$, the weight decay and the gradient-clipping norm matched as well. QuRL's ablation raises the global learning rate instead, so it does not control for this.
\end{itemize}

We confirm both results on real stacks. With vLLM's INT8 kernel, UAQ leaves the rollouts of 0.5B--7B models bit-identical at $s=2$, and leaves their rollout--trainer mismatch unchanged at $s=1.5$. A full QuRL-style training loop with the matched settings reproduces UAQ bit for bit over 80 updates.

We also show that re-quantization does not freeze the rollout policy: a sub-grid update flips a code with probability $|\delta|/\Delta$. What re-quantization does create is a rounding error that stays constant in size but persists over training. A pre-registered, instrumented PPO run confirms both on the synced master weights. From the second update on, 97--99\% of each update is captured, and after 435 steps the rounding error still correlates 0.69 with its initial value, as predicted. The first update from the BF16 checkpoint is captured at only 93\%, because BF16 storage breaks the theory's uniform-phase assumption at the scale of one update.

Finally, QuRL's stated GSM8K settings are those of VeRL's PPO quickstart, whose actor learning rate is $10^{-6}$. At the $10^{-5}$ that QuRL states, our runs did not learn; at $10^{-6}$, BF16 rollouts match the published baselines. In a pre-registered six-arm experiment at that rate, UAQ left the rollout--trainer mismatch unchanged relative to a BF16 control, while its effects on accuracy were too small to resolve with three to five seeds per arm.
\end{abstract}
